% Baccalauréat malgache, série D, 2015 : énoncé.
\TCDBacAttribution{D}{2015}
\noindent \underline{\textsc{Exercice I}} (5 points) \newline
      Soit $f$ une transformation d\'efinie dans un plan complexe $(P)$ par :
      \[
      \begin{array}[t]{rcl}
            f : P & \rightarrow & P \\
             M(z) & \rightarrow & M'(z')
      \end{array}
      \]
      telle que 
      $f(z) = z' = \frac{iz - 5 + i}{z-4}$.
      \begin{enumerate}
            \item \begin{enumerate}
                        \item D\'eterminer l'ensemble de d\'efinition de $f$ not\'e $Df$.
                        \item R\'esoudre dans $\mathbb{C}$ l'\'equation $f(z) = z$.
                  \end{enumerate}
            \item On consid\`ere dans le plan complexe $(P)$ muni d'un rep\`ere orthonorm\'e $(O;\vec{u},\vec{v})$
                  d'unit\'e 1~cm les points $A$, $B$, $C$ et $M$ d'affixes respectives 
                  $-2+i$, $1-i$, $3+2i$ et $z = x + i y$ o\`u $x,\ y \in \mathbb{R}$.
                  \begin{enumerate}
                        \item D\'eterminer et construire l'ensemble $(\mathscr{D})$ des 
                              points $M(z)$ pour que $f(z)$ soit r\'eel.
                        \item On pose $Z = \frac{z_A - z_B}{z_C - z_B}$.
                              Donner la forme trigonom\'etrique de $Z$ et en d\'eduire la nature du triangle $ABC$.
                  \end{enumerate}
            \item Calculer l'affixe du point $D$ tel que le quadrilat\`ere $ABCD$ soit un carr\'e.
            \item Soit $S$ la similitude plane directe de rapport $\sqrt{2}$, d'angle $\frac{\pi}{4}$ et $S(B) = C$,
                  dont son \'ecriture complexe est de la forme $S : z' = a z + b$ o\`u $a\in\mathbb{C}^*$ et $b\in\mathbb{C}$.
                  \begin{enumerate}
                        \item Donner la forme alg\'ebrique du nombre complexe $a$.
	                        \item En d\'eduire la valeur du nombre complexe $b$.
	                        \item Pr\'eciser le centre de $S$.
                  \end{enumerate}
      \end{enumerate}
%%%%%%%%%%%%%%%%%%%%%%%%%%%%%%%%%%%%%%%%%%%%%%%%%%%%%%%%%%%%%%%%%%%%%%%%%%%%%%%%%%%%%%%%%%%%%%%
%%%%%%%%%%%%%%%%%%%                   %%%%%%%%%%%%%%%%%%%                   %%%%%%%%%%%%%%%%%%%
%%%%%%%%%%%%%%%%%%%%%%%%%%%%%%%%%%%%%%%%%%%%%%%%%%%%%%%%%%%%%%%%%%%%%%%%%%%%%%%%%%%%%%%%%%%%%%%
\noindent \underline{\textsc{Exercice II}} (5 points) 
      \begin{description}
            \item[I.]
	                  Douze jetons indiscernables au toucher, plac\'es dans une urne, portent les lettres : 
                  A, A, A, A, B, B, B, D, D, E, E, E.
                  \begin{enumerate}
                        \item On tire successivement, avec remise, trois jetons de l'urne.
	                              Calculer les probabilit\'es des \'ev\`enements suivants : 
                              \begin{description}
                                    \item[G :] << Obtenir au moins une lettre B >>.
                                    \item[H :] << Obtenir exactement une voyelle >>.
                              \end{description}
	                        \item On tire simultan\'ement quatre jetons de l'urne.
                              On d\'esigne par $X$ la variable al\'eatoire associ\'ee au nombre de consonnes obtenues.
                              \begin{enumerate}
                                    \item D\'eterminer l'univers image de $X$.
                                    \item Donner la loi de probabilit\'e de $X$.
                                    \item D\'efinir et repr\'esenter graphiquement la loi de probabilit\'e de $X$.
                              \end{enumerate}
                  \end{enumerate}
	                  \underline{NB :} On donnera les r\'esultats sous forme de fraction irr\'eductible.
            \item[II.]
                  Les chiffres d'affaires d'une entreprise de l'ann\'ee 2008 \`a 2012 
                  sont repr\'esent\'es dans le tableau suivant : \newline
                  \begin{tabular}[b]{|p{20ex}|c|c|c|c|c|}
                  \hline 
                  Ann\'ee & 2008 & 2009 & 2010 & 2011 & 2012 \\ 
                  \hline 
                  Rang de l'ann\'ee $x_i$ & 0 & 1 & 2 & 3 & 4 \\ 
                  \hline 
                  Chiffre d'affaire en millions d'ariary & 504 & 580 & 644 & $y_3$ & 735 \\ 
                  \hline 
                  \end{tabular}.
                  
                  L'\'equation de la droite de r\'egression $(D)$ de $y$ en $x$ est : $y = 57,3 x + 516,2$.
                  \begin{enumerate}
                        \item Calculer les coordonn\'ees du point moyen $G$.
                        \item En d\'eduire la valeur de $y_3$.
                        \item En quelle ann\'ee l'entreprise pourra-t-elle atteindre le chiffre d'affaire 
                              d'un milliard quatre cent trente trois millions d'Ariary\,?
                  \end{enumerate}
      \end{description}
%%%%%%%%%%%%%%%%%%%%%%%%%%%%%%%%%%%%%%%%%%%%%%%%%%%%%%%%%%%%%%%%%%%%%%%%%%%%%%%%%%%%%%%%%%%%%%%
%%%%%%%%%%%%%%%%%%%                   %%%%%%%%%%%%%%%%%%%                   %%%%%%%%%%%%%%%%%%%
%%%%%%%%%%%%%%%%%%%%%%%%%%%%%%%%%%%%%%%%%%%%%%%%%%%%%%%%%%%%%%%%%%%%%%%%%%%%%%%%%%%%%%%%%%%%%%%
\noindent \underline{\textsc{Probleme}} \newline
      Soit $f$ la fonction num\'erique d\'efinie sur $Df = ]-1 ; +\infty[$ par : $f(x) = \ln (x+1) + \expo{-x}$.
      On note $(\mathscr{C})$ la courbe repr\'esentative dans un rep\`ere orthonorm\'e direct $(O;\ivec, \jvec)$
      d'unit\'e graphique 1~cm.
      \begin{enumerate}
            \item Soit $g$ la fonction num\'erique d\'efinie sur $\mathbb{R}$ par : $g(x) = x+1 - \expo{x}$.
                  \begin{enumerate}
                        \item Etudier les variations de $g$ 
                              (on ne demande pas de calculer les limites en $-\infty$ et $+\infty$).
                        \item En d\'eduire le signe de $g(x)$ suivant les valeurs de $x$.
                  \end{enumerate}
            \item \begin{enumerate}
                        \item Calculer $\displaystyle{\lim_{x\rightarrow -1^+} f(x)}$.
                              Interpr\'eter graphiquement ce r\'esultat.
                        \item Calculer $\displaystyle{\lim_{x\rightarrow +\infty} f(x)}$.
                        \item Montrer que pour tout $x\in Df$ :
                              \[
                                    \frac{f(x)}{x} = \frac{\ln (x+1)}{x+1} \times \frac{x+1}{x} + \frac{1}{x \expo{x}} \;.
                              \]
                        \item En d\'eduire $\displaystyle{\lim_{x\rightarrow +\infty} \frac{f(x)}{x}}$.
                              Que peut-on conclure sur la courbe $(\mathscr{C})$\,?
                  \end{enumerate} 
            \item \begin{enumerate}
                        \item Exprimer $f'(x)$ en fonction de $g(x)$
                              pour tout $x\in Df$ o\`u $f'$ est la fonction d\'eriv\'ee de $f$.
                        \item Dresser le tableau de variation de $f$ sur $Df$.
                  \end{enumerate}
            \item \begin{enumerate}
                        \item Montrer que $(\mathscr{C})$ coupe l'axe des abscisses en un point unique 
                              d'abscisse $\alpha \in ]-1 ; 0 [$.
                        \item Donner une \'equation de la tangente $(T)$ \`a $(\mathscr{C})$ 
                              au point d'abscisse $x_0 = 0$.
                  \end{enumerate}
            \item Construire $(\mathscr{C})$ et $(T)$ dans le m\^eme rep\`ere.
            \item \begin{enumerate}
                        \item D\'eterminer les r\'eels :
                              \[
                                    \frac{x}{x+1} = a + \frac{b}{x+1}       \;.
                              \]
                        \item \`A l'aide d'une integration par partie, calculer une primitive de la fonction 
                              $k : x \mapsto \ln (x+1)$ sur $Df$.
                        \item Calculer en $\mathrm{cm}^2$ l'aire g\'eom\'etrique du domaine plan 
                              limit\'e par la courbe $(\mathscr{C})$, l'axe des abscisses
                              et les droites d'\'equations $x=0$, $x=1$.
                  \end{enumerate}
            \item \begin{enumerate}
                        \item Montrer que $f$ r\'ealise une bijection de $Df$ 
                              vers un intervalle $J$ que l'on d\'eterminera.
                        \item Construire la courbe $(\mathscr{C}')$ qui est la repr\'esentation 
                              graphique de $f^{-1}$, fonction r\'eciproque de $f$ dan le m\^eme 
                              rep\`ere que $(\mathscr{C})$.
                  \end{enumerate}
      \end{enumerate}









